% CHECKPOINT35: complete formalization of the frozen checkpoint34 residual proof.
% Mathematical predecessor SHA256 e34097f6bebd644c939c5aa4bef1bc13e854695ccc8430d30a640db6fa787422.
\statement{A.116}{Proved}{The relative-perfect congruence residual and its full quotient action}
Retain the original negative full-\(\widehat{\mathrm{GT}}\) involution and
its actual canonical data from A.103--A.108. Let
\[
 \begin{gathered}
 B=C_2(s)\amalg C_3(r),\quad F=\ker(B\to S_3),\quad
 N=\ker(q:F\to C),\quad N_1=\overline{[N,F]},\\
 P=F^{(2)},\quad \pi_2:F\to P,\quad O=\ker\pi_2,\quad
 R=N_1\cap O=N\cap F'\cap O.
 \end{gathered}                                                   \tag{A.116.1}
\]
Put \(C_o=\prod_{p\text{ odd}}\mathrm{SL}_2(\mathbb Z_p)\) and
\(D_o=(F_{\rm ab})_o\times_{(C_o)_{\rm ab}}C_o\), using the actual
abelianization maps. Then
\[
 \boxed{F/R=P\times D_o,\qquad R=\overline{[R,B]},\qquad H_d\in R.}
                                                                  \tag{A.116.2}
\]
The original canonical substitution \(\psi_d\) induces an automorphism
of \(F/R\). The assignments \(s\mapsto v=u^{-1}s\),
\(\delta=sr\mapsto\delta^d\) define an actual automorphism
\(\Psi_R\) of \(B/R\), with \(\alpha\Psi_R=\Psi_Rc_B\).
There exists a full peripheral automorphism \(\Gamma_R\) of \(B\)
inducing this specified \(\Psi_R\) and satisfying
\(\Gamma_R(\delta)=\delta^d\) literally.
The group \(R\) is not pro-two and is nontrivial. In the notation of
A.115, \(1\ne R\le L_\infty\); that central tower is therefore not
cofinal. No assertion \(R=R'\), \(R=L_\infty\), or \(H_d=1\) is made.

\begin{proof}
\emph{1. The actual prior coordinates and stability.}
A.106 gives
\[
 N_1=N\cap F'\cap\pi_2^{-1}(M_{21}),\quad
 M_2=\ker(P\to C_{(2)}),\quad M_{21}=\overline{[M_2,P]}.
                                                                  \tag{A.116.3}
\]
Intersecting with \(O\) proves the equality defining \(R\).
The geometric signed parity argument of A.107 makes \(N,N_1\)
normal in \(B\), and they are stable under the original \(\alpha\)
and \(c_B\). The group \(O\) is characteristic in \(F\), and
\(\psi_d\) preserves it by functoriality of the maximal pro-two
quotient. A.104/A.107 prove the other required \(\psi_d\)-stabilities.
Thus all quotient maps used here are actual continuous maps; no
characteristic-kernel assertion for arbitrary automorphisms of \(F\)
is being substituted for geometric stability.

A.114 proves, with \(\overline B=B/O\) and \(B_1=B/N_1\),
\[
 \begin{gathered}
 \pi_2(N_1)=M_{21},\qquad
 H^2(B_1,\mathbb Q_\ell/\mathbb Z_\ell)=0\quad(\ell\text{ odd}),\\
 H^1(D_o,\mathbb Q_2/\mathbb Z_2)
 =H^2(D_o,\mathbb Q_2/\mathbb Z_2)=0,\qquad
 H^2(\overline B,\mathbb Q_2/\mathbb Z_2)=0.
 \end{gathered}                                                    \tag{A.116.4}
\]
These are the specific continuous cohomology groups calculated there,
using Beyl's Proposition~5.7 and the complete product and prime-three
lattice arguments. There is no vanishing assumption about
\(H^2(C_{(2)},\mathbb Q/\mathbb Z)\).

\medskip\noindent\emph{2. The full-pro-two product.}
Consider the marked continuous map
\[
 F\longrightarrow P\times D_o,\qquad
 g\longmapsto\bigl(\pi_2(g),(g_{{\rm ab},o},q_o(g))\bigr).
                                                                  \tag{A.116.5}
\]
Its kernel is \(R\): the first coordinate already forces the dyadic
signed and dyadic abelian coordinates to vanish, and the other two
then vanish exactly when \(g\in N\cap F'\).
Its \(D_o\)-projection is onto. Given a compatible pair \((a,c)\),
choose an element with the prescribed odd abelian vector. Its
remaining signed discrepancy is in \(C_o'\), and is corrected by an
\(F'\)-lift, since continuous compact surjections map the closed
derived subgroup onto the closed derived subgroup.
The \(P\)-projection is onto by definition.

The odd abelianization \((C_o)_{\rm ab}\) is \(C_3\), supported at
three, and the original character is the sum of the two coordinates
modulo three. The kernel of \(D_o\to C_o\) is pro-odd abelian.
Every finite two-group quotient kills that kernel. A nontrivial
two-group quotient of \(C_o\) would admit a nontrivial \(C_2\)
character, contradicting its abelianization. Hence \(D_o\) has no
nontrivial pro-two quotient. Profinite Goursat applied to the closed
subdirect image of (A.116.5) gives the actual isomorphism
\(F/R=P\times D_o\). In particular the image of \(O\) is exactly
\(D_o\), which is normal in \(B/R\), and
\[
 (B/R)/D_o=\overline B,\qquad
 \ker(B/R\to B_1)=N_1/R=M_{21}.
                                                                  \tag{A.116.6}
\]
The latter kernel is pro-two; it is not asserted central.

\medskip\noindent\emph{3. Relative perfectness.}
For an odd prime \(\ell\), the pro-two kernel in (A.116.6) has zero
positive cohomology with the trivial discrete module
\(\mathcal A_\ell=\mathbb Q_\ell/\mathbb Z_\ell\), by finite-quotient
averaging. The actual extension spectral sequence and (A.116.4) give
\(H^2(B/R,\mathcal A_\ell)=H^2(B_1,\mathcal A_\ell)=0\).
For two-primary coefficients, the first extension in (A.116.6)
and the two vanishings for \(D_o\) in (A.116.4) give the inflation
isomorphism
\[
 H^2(\overline B,\mathcal A_2)\xrightarrow{\sim}
 H^2(B/R,\mathcal A_2)=0.
\]
No higher two-primary cohomology vanishing for \(D_o\) is required.
Continuous cochains into the discrete torsion group \(\mathbb Q/\mathbb Z\)
have finite image; splitting their finitely many primary components
therefore proves
\(H^2(B/R,\mathbb Q/\mathbb Z)=0\).

In the five-term sequence for \(R\to B\to B/R\), restrictions of
characters of \(B\) to \(R\) vanish because \(R\subset F'\subset B'\).
The degree-two vanishing consequently gives
\(H^1(R,\mathbb Q/\mathbb Z)^B=0\).
These are exactly the characters of the compact abelian quotient
\(R/\overline{[R,B]}\), and continuous characters separate its points.
Thus \(R=\overline{[R,B]}\), which is relative perfectness in \(B\).

The same original word \(H_d\) belongs to \(N_1\) by A.106, and its
entire pro-two image is one by A.89.3, using A.42/A.45 for this actual
\(u,d\). Therefore \(H_d\in N_1\cap O=R\).
This uses full word vanishing, not merely a dyadic coefficient or a
finite collection of congruence heads.

\medskip\noindent\emph{4. The canonical quotient automorphisms.}
In (A.116.5) the induced \(\psi_d\)-action is the product of its
actual endomorphism of \(P\) and
\[
 (a,c)\longmapsto\bigl(da,\operatorname{Ad}(H)c\bigr)
 \quad\text{on }D_o,\qquad
 H=\begin{pmatrix}1&-b\\0&d^{-1}\end{pmatrix}.
                                                                  \tag{A.116.7}
\]
Its action on the Frattini quotient of the finitely generated
free pro-two \(P\) is \(dI\), so the pro-two Frattini criterion and
Hopficity make the first map an automorphism. In the second map,
both coordinate maps are automorphisms. Their common \(C_3\)
action is multiplication by the unit \(d\bmod3\), by the actual
abelianization identity of \(\psi_d\); their inverses consequently
preserve the fiber-product condition too. No cross term is missing:
all coordinates in (A.116.5) are the literal original quotient maps.
Thus \(\psi_d\) is an automorphism on \(F/R\).

The original ordered identity
\[
 H_d=u\bigl((v\delta^d)^3\bigr)^{-1}u^{-1},\qquad v=u^{-1}s,
                                                                  \tag{A.116.8}
\]
now proves the cubic before a modular map is asserted. Together with
\(v^2=1\), it defines \(B\to B/R\) by
\(s\mapsto v\), \(r\mapsto v\delta^d\).
Its pure restriction is the actual \(\psi_d\), so it kills \(R\)
and descends. It is an automorphism on \(F/R\) and the identity on
\(S_3\), hence an automorphism \(\Psi_R\) of \(B/R\).
The original generator identities give
\(\alpha\Psi_R=\Psi_Rc_B\) there.

\medskip\noindent\emph{5. Lifting precisely this quotient action.}
The involution \(v\) already has its full factor frame in \(B\), by
finite-factor torsion conjugacy and the original \(S_3\) normalization.
The A.108 order-three frame lifts from \(B_1\) through the pro-two
kernel \(M_{21}\) in (A.116.6): the residual \(C_3\)-cocycle is
killed by the complete continuous coprime lemma in A.108.
Thus \(\Psi_R\) has the actual factor frames with the exponent
\(e=d\bmod3\), and its cusp image is literally \(\delta^d\).

Choose an actual Galois peripheral automorphism \(\gamma\) of
multiplier \(d\). The original \(\mathrm{LM7\_PUSH}\) E.65.8 and
E.65.15 cusp and signed covariance formulas, in the exact based
marking used in A.114, show that \(\gamma(N)=N\).
It preserves \(F'\), \(O\), and therefore \(R\).
Left composition by \(\gamma^{-1}\) turns \(\Psi_R\) into a
positive-peripheral automorphism of \(B/R\); the unit powers have
now actually been removed. The closed-normal modular criterion
of A.3/A.34 requires a positive-peripheral lift to
\(B/\overline{[R,B]}\). By relative perfectness this is the same
quotient \(B/R\), on which the required map is already defined.
The criterion yields a full positive-peripheral lift with the specified
quotient frames. Its cusp conjugator lies in \(R\); correcting by it
fixes the cusp literally without changing the action on \(B/R\).
Composing \(\gamma\) back gives \(\Gamma_R\) with the asserted exact
quotient action and cusp value.

This conclusion does not identify the pure restriction of
\(\Gamma_R\) with \(\psi_d\) on the full \(F\), nor assert
\(\alpha\Gamma_R=\Gamma_Rc_B\) on the full \(B\).

\medskip\noindent\emph{6. The residual is nontrivial.}
A.110 proves that the actual \(N_1\) is not pro-two.
Since \(N_1/R=M_{21}\) is pro-two, \(R\) cannot be pro-two: otherwise
\(N_1\) would be an extension of pro-two groups. In particular
\(R\ne1\). Relative perfectness gives \(R\le L_j\) at every
stage of A.115 by induction, hence \(1\ne R\le L_\infty\).
Equality with that intersection would require a separate calculation
in \(\overline B\), and is not inferred here.
\end{proof}
\noindent\emph{Scope.}
All groups and closures are the actual ambient profinite objects.
The original canonical word remains unevaluated inside this nontrivial
residual. Both main problems and the coherent CM obstruction remain open.
